\DTMsavedate{mydate}{2022-06-30}
\maketitle{The Professional Engineer \Romannum{1}}{Electronic Engineering}{\DTMusedate{mydate}}

% ----------------------------------------------------- %
% COURSE INFO
\subsection*{Course Information}
\vspace{0ex}%
\noindent\parbox{0.5\textwidth}{%
    \noindent\begin{tabular}{@{}ll}
                 \textsf{Course:}          & 	The Professional Engineer \Romannum{1} \\
                 \textsf{Course Code:}         & 	TEE 2205    \\
                 \textsf{Credits:}    & 	10
    \end{tabular}}
\vspace{2ex}\hrule\vspace{2ex}

% ----------------------------------------------------- %
% INSTRUCTOR INFO
\subsection*{Lecturer Information}
\noindent\parbox{0.5\textwidth}{%
    \noindent\begin{tabular}{@{}ll}
                 \textsf{Name:}         &    Fidelis Nhenga Mugarisanwa \\
% \textsf{Phone:}        &    \phone          \\
\textsf{Email:}        &    fidelis.mugarisanwa@nust.ac.zw \\
\textsf{Qualifications:}        & BSc in Computer Science
    \end{tabular}}
\vspace{2ex}\hrule\vspace{2ex}


% ----------------------------------------------------- %
% COURSE SYNOPSIS
\subsection*{Course Synopsis}
The course aims at promoting the student’s understanding of his/her role in society as the future engineer.
At the end of the course, the student should be able to at least appreciate his or her duty as a potential innovator and/or service provider who must act within the moral-legal framework.
The student must also be able to link every engineering activity with the needs of society and be conscious of everyday problems that may challenge their ability to generate solutions to some of those problems.
Every possible attempt will be made to encourage the student to keep abreast with ever-changing technology and the subsequent demands imposed on him or her from the community that they live in.
The students should also be able to independently manage or at least develop interest in small-scale multidisciplinary projects.
\vspace{2ex} \hrule\vspace{2ex}

% ----------------------------------------------------- %
% COURSE TOPICS
\subsection*{Topics}
\begin{outline}
    \1 Introduction
    \1 Idea Generation and Product Development
    \1 Engineering Business Environment
    \1 Business Entity Types
    \1 Business Models
    \1 Research Methods
    \1 Project Management
    \1 Ethics in Engineering
    \1 Safety in the Workplace
\end{outline}
\vspace{2ex} \hrule\vspace{2ex}

% ----------------------------------------------------- %
% Students Assenssment
\subsection*{Assessment of Students}
\begin{outline}
    \1 Continuous assessment will consist of in-class presentations, written assignments and written in-class tests.
    \1 There shall be a final examination at the end of the semester.
\end{outline}
\vspace{2ex}\hrule\vspace{2ex}

% ----------------------------------------------------- %
% Grade Evaluation
\subsection*{Course Evaluation and Grading Policies}
Grades are based on:
Continuous assessment (\qty{25}{\percent}),
Final Exam (\qty{75}{\percent})
\vspace{2ex}\hrule
\vspace{2ex}\hrule\vspace{2ex}